\documentclass{article}

% --- Fermion TeX Engine demo document -------------------------------
% Macros are tracked by the dependency graph: change a definition below
% and only the blocks that actually use it are re-expanded.
\newcommand{\term}[1]{\textbf{#1}}
\newcommand{\engine}{Fermion TeX Engine}
\newcommand{\half}{\frac{1}{2}}

\begin{document}

\section{Introduction}
\label{sec:intro}

This document is being typeset by the \engine, a resident incremental
TeX runtime. Unlike a batch compiler, the engine keeps the entire
document state alive between keystrokes: the \term{Source DOM}, the
\term{Semantic DOM}, the dependency graph, the \term{Layout DOM}, the
\term{Page DOM} and the display lists all stay in memory while you
edit. Nothing is re-read from disk and no process is restarted.

When you change a single word of this paragraph, the engine marks
exactly one source node dirty, re-expands and re-parses only that
block, re-runs line breaking for one paragraph, and patches only the
pages whose display list actually changed. Try it now: edit any word
and watch the inspector on the right report the dirty set.

The core claim of this project is simple to state. Given an edit, the
engine must answer one question precisely: what becomes dirty? The
answer flows through every layer, from the changed source range down
to the display commands of a single page, and the inspector shows the
whole chain for every keystroke.

\section{The DOM Pipeline}
\label{sec:pipeline}

As explained in Section~\ref{sec:intro}, the engine is not an editor
and not a wrapper around an existing TeX binary. It is a runtime. The
source text is segmented into blocks, each block is expanded by the
macro VM, parsed into a semantic node, laid out into lines, and the
lines are poured into pages. Every stage caches its result per block
and every cache records what it depends on.

The pipeline layers are:

\begin{itemize}
\item Source DOM: blocks with stable identity across edits, diffed by
  content hash so an unchanged paragraph keeps its caches.
\item Macro VM: expansion with per-block dependency records, so
  redefining a macro invalidates exactly its users.
\item Semantic DOM: sections, paragraphs, lists, math, labels, refs.
\item Dependency graph: labels, references and counters, used to
  propagate dirtiness beyond the edited block.
\item Layout DOM: justified line breaking with real font metrics.
\item Page DOM: incremental pagination with convergence detection.
\item Display lists: absolute drawing commands, patched per page.
\end{itemize}

Inline mathematics is measured and placed by the engine itself, for
example $E = mc^2$ or $\alpha + \beta \le \gamma$, and participates in
line breaking as a rigid box. Display equations are blocks of their
own:

\[
a^2 + b^2 = c^2
\]

\[
f(x) = \half x^2 + \alpha x + \beta
\]

Because the equation above uses the macro \texttt{half}, redefining
that macro in the preamble re-typesets exactly one block. The
dependency graph records the relationship the moment the expansion
happens, so no global scan is ever needed.

\section{Incremental Pagination}
\label{sec:pages}

Page breaking is the hardest part of incremental typesetting. A small
change of paragraph height can shift every following page. The engine
therefore treats pages as chunks of a stream of line units, where each
unit is owned by the layout node of its block. Unchanged blocks
contribute the very same unit objects as in the previous pass, so two
runs that start a page at the same unit will break all following
pages identically.

This gives convergence for free: after an edit, pagination is
recomputed from the stream, but every page whose unit sequence is
identical to the previous run is adopted wholesale, display list and
all. The inspector reports how many pages were reused and how many
were rebuilt, and the patch set contains only pages whose display
list hash actually changed.

Cross references stress the dependency graph in the other direction.
This sentence refers to Section~\ref{sec:pipeline} and to
Section~\ref{sec:intro}; renaming either label makes exactly the
blocks containing these references dirty, because the graph knows
which semantic nodes consumed which labels. Deleting a label turns
the reference into a visible question mark pair, just as in LaTeX.

\subsection{Convergence in practice}
\label{sec:convergence}

Try inserting a full paragraph near the top of the document. Every
following line shifts, so the pages after the edit are rebuilt, and
the inspector shows the pagination cost honestly. Then undo the
insertion: the stream returns to its previous shape, the reused page
count jumps back up, and the patch set shrinks to the pages that
actually changed on screen.

Now try editing this very paragraph instead. The blocks before it are
untouched, their layout nodes are cache hits, and the pages before
this one are adopted without regeneration. Only the page containing
this paragraph, and possibly the one after it, appear in the patch
list. That asymmetry between where you type and what gets recomputed
is the whole point of a resident typesetting runtime.

\subsection{What the engine refuses to fake}
\label{sec:honest}

The preview in this demo is not a bitmap diff and not a reloaded PDF.
The viewer receives display lists, one per changed page, and redraws
exactly those pages. The PDF export button serializes the same
display lists into a real PDF file, which demonstrates that the
document state is upstream of any particular output format.

The engine also refuses to hide its limits. Unknown commands are
rendered visibly in typewriter type and reported as diagnostics
instead of being silently dropped. Unclosed groups are tolerated so
the runtime survives every intermediate keystroke while you type a
command name character by character.

\section{Numbers and Counters}
\label{sec:counters}

Section numbers are computed by the counter pass, and headings whose
number changed are re-laid-out even when their text did not change.
Inserting a new section before this one renumbers everything that
follows, and the dependency report attributes those relayouts to the
counter, not to the edit itself. References such as the one pointing
at Section~\ref{sec:convergence} update in the same pass.

The demo ends with a short summary. The engine keeps a living DOM of
the document, tracks dependencies between macros, labels, counters
and blocks, lays out only what changed, repaginates only from where
heights changed, and patches only pages whose drawing commands
differ. Everything else is reuse, and the inspector proves it on
every keystroke.

\end{document}
